% Options for packages loaded elsewhere
\PassOptionsToPackage{unicode}{hyperref}
\PassOptionsToPackage{hyphens}{url}
\PassOptionsToPackage{space}{xeCJK}
\documentclass[
  chinese,
  11pt,
]{ctexart}
\usepackage{xcolor}
\usepackage[margin=2.2cm]{geometry}
\usepackage{amsmath,amssymb}
\setcounter{secnumdepth}{5}
\usepackage{iftex}
\ifPDFTeX
  \usepackage[T1]{fontenc}
  \usepackage[utf8]{inputenc}
  \usepackage{textcomp} % provide euro and other symbols
\else % if luatex or xetex
  \usepackage{unicode-math} % this also loads fontspec
  \defaultfontfeatures{Scale=MatchLowercase}
  \defaultfontfeatures[\rmfamily]{Ligatures=TeX,Scale=1}
\fi
\usepackage{lmodern}
\ifPDFTeX\else
  % xetex/luatex font selection
  \ifXeTeX
    \usepackage{xeCJK}
    \setCJKmainfont[]{Noto Serif SC}
    \setCJKsansfont[]{Noto Sans SC}
    \setCJKmonofont[]{DejaVu Sans Mono}
  \fi
  \ifLuaTeX
    \usepackage[]{luatexja-fontspec}
    \setmainjfont[]{Noto Serif SC}
  \fi
\fi
% Use upquote if available, for straight quotes in verbatim environments
\IfFileExists{upquote.sty}{\usepackage{upquote}}{}
\IfFileExists{microtype.sty}{% use microtype if available
  \usepackage[]{microtype}
  \UseMicrotypeSet[protrusion]{basicmath} % disable protrusion for tt fonts
}{}
\makeatletter
\@ifundefined{KOMAClassName}{% if non-KOMA class
  \IfFileExists{parskip.sty}{%
    \usepackage{parskip}
  }{% else
    \setlength{\parindent}{0pt}
    \setlength{\parskip}{6pt plus 2pt minus 1pt}}
}{% if KOMA class
  \KOMAoptions{parskip=half}}
\makeatother
\usepackage{longtable,booktabs,array}
\newcounter{none} % for unnumbered tables
\usepackage{calc} % for calculating minipage widths
% Correct order of tables after \paragraph or \subparagraph
\usepackage{etoolbox}
\makeatletter
\patchcmd\longtable{\par}{\if@noskipsec\mbox{}\fi\par}{}{}
\makeatother
% Allow footnotes in longtable head/foot
\IfFileExists{footnotehyper.sty}{\usepackage{footnotehyper}}{\usepackage{footnote}}
\makesavenoteenv{longtable}
\ifLuaTeX
\usepackage[bidi=basic,shorthands=off,provide=*]{babel}
\else
\usepackage[bidi=default,shorthands=off,provide=*]{babel}
\fi
\ifLuaTeX
  \usepackage{selnolig} % disable illegal ligatures
\fi
\setlength{\emergencystretch}{3em} % prevent overfull lines
\providecommand{\tightlist}{%
  \setlength{\itemsep}{0pt}\setlength{\parskip}{0pt}}
\usepackage{bookmark}
\IfFileExists{xurl.sty}{\usepackage{xurl}}{} % add URL line breaks if available
\urlstyle{same}
\hypersetup{
  pdftitle={课题06-多卡与分布式},
  pdflang={zh-CN},
  hidelinks,
  pdfcreator={LaTeX via pandoc}}

\title{课题06-多卡与分布式}
\author{}
\date{}

\begin{document}
\maketitle

{
\setcounter{tocdepth}{2}
\tableofcontents
}
\section{课题06 开题简报 · 多卡与分布式训练}\label{ux8bfeux989806-ux5f00ux9898ux7b80ux62a5-ux591aux5361ux4e0eux5206ux5e03ux5f0fux8badux7ec3}

\begin{quote}
模块：\textbf{M4} ｜ 周次：\textbf{W9（11-09→11-15）} ｜ 算力：\textbf{T2 Kaggle 2×T4，计划 8 GPU·小时} 唐杰作业对应：\textbf{② 的''多卡训练增益''部分} ｜ 判分来源：\textbf{CS336 A2 分布式部分 + 自建 benchmark} 手写：\textbf{R3 的延伸（并行策略选择与通信原语判断）}
\end{quote}

\subsection{一、研究背景}\label{ux4e00ux7814ux7a76ux80ccux666f}

数据并行（DDP）把同一模型复制到每张卡、各自算一个微批、再\textbf{同步梯度}； 当模型放不下单卡时，才有 FSDP/ZeRO（切分参数、梯度、优化器状态）、张量并行（切层内矩阵）、流水线并行（切层）。 \textbf{Kaggle 免费档给的是 2×T4，且走 PCIe 而非 NVLink} ------ 这恰好是''观察通信成为瓶颈''的最佳实验台。

\subsection{二、研究问题}\label{ux4e8cux7814ux7a76ux95eeux9898}

\begin{quote}
\textbf{在弱互联（PCIe）的 2×T4 上，DDP 的扩展效率是多少？瓶颈是计算还是通信？FSDP 在什么规模才开始划算？}
\end{quote}

\subsection{三、攻关线索}\label{ux4e09ux653bux5173ux7ebfux7d22}

\begin{enumerate}
\def\labelenumi{\arabic{enumi}.}
\tightlist
\item
  \textbf{先对齐正确性}：DDP 后的梯度应与单卡（同等效批大小）\textbf{数值一致}，先验证再谈性能
\item
  \textbf{批大小语义}：2 卡 × 每卡 batch = 等效批大小；若显存受限要\textbf{同步调整学习率}（线性/平方根缩放）
\item
  \textbf{计时口径}：\texttt{step\ time} 要区分''纯计算时间''与''含 all-reduce 的时间''；用 \texttt{torch.profiler} 或手写计时
\item
  \textbf{扩展效率}：\(E = \dfrac{T_1}{2\,T_2}\)。\textbf{先预测}（PCIe 下 0.6--0.8？）再测
\item
  \textbf{通信占比}：梯度同步量与参数量成正比（\(2N\) 字节/步，all-reduce 近似）；序列长/批小时通信占比更高
\item
  \textbf{FSDP 的代价}：切分带来额外通信（all-gather + reduce-scatter），\textbf{小模型上往往更慢}
\end{enumerate}

\subsection{四、里程碑}\label{ux56dbux91ccux7a0bux7891}

\begin{itemize}
\tightlist
\item[$\square$]
  M1 DDP 梯度一致性验证（与单卡同等效批）→ 记录最大误差
\item[$\square$]
  M2 单卡 vs 双卡 的 \texttt{step\ time}、\texttt{tokens/s} 对照（同模型、同等效批）
\item[$\square$]
  M3 \textbf{扩展效率曲线}（至少 3 个配置点：不同 batch/序列长度）
\item[$\square$]
  M4 \textbf{通信占比分解}：计算 / all-reduce / 数据加载 三段计时
\item[$\square$]
  M5 一次 FSDP 实验并解释''为什么（不）划算''
\item[$\square$]
  M6 \textbf{先写预测再实测}：预测扩展效率，跑完对照（进 \texttt{memory/MEMORY.md} 打脸日志）
\item[$\square$]
  M7 一页报告
\end{itemize}

\subsection{五、交付物}\label{ux4e94ux4ea4ux4ed8ux7269}

{\def\LTcaptype{none} % do not increment counter
\begin{longtable}[]{@{}lll@{}}
\toprule\noalign{}
交付物 & 位置 & 验收 \\
\midrule\noalign{}
\endhead
\bottomrule\noalign{}
\endlastfoot
分布式脚本 & \texttt{手写/} 或 \texttt{脚手架/}（按 B3 分类） & 可复跑 \\
计时原始数据 & \texttt{证据/bench.csv} & 含每配置 3 次重复 \\
扩展效率曲线 & \texttt{证据/} & 标注硬件（2×T4 PCIe）与对比基准 \\
报告 & \texttt{报告.md} & 解释''为什么不是 2 倍'' \\
\end{longtable}
}

\subsection{六、讲课锚点}\label{ux516dux8bb2ux8bfeux951aux70b9}

\begin{itemize}
\tightlist
\item
  \textbf{讲义}：\texttt{M4-训练.md}（W3 展开）中并行一节
\item
  \textbf{对标}：CS336 lecture 07--08（Parallelism）· \texttt{nanochat/optim.py}（分布式优化器实现）· HF Ultra-Scale Playbook（用前先核）
\item
  \textbf{搭桥}：M0 的''梯度累加''→ DDP 就是\textbf{跨卡的梯度累加}
\end{itemize}

\subsection{七、代码级提问示例}\label{ux4e03ux4ee3ux7801ux7ea7ux63d0ux95eeux793aux4f8b}

\begin{enumerate}
\def\labelenumi{\arabic{enumi}.}
\tightlist
\item
  2×T4 上 DDP 只快了 1.4 倍，列出三个最可能的原因。
\item
  把每卡 batch 减半、用梯度累积补回等效批，通信量会怎么变？
\item
  什么情况下 FSDP 会\textbf{比 DDP 慢}？（用通信量解释）
\item
  all-reduce 的通信量是多少字节/步？用参数量 \(N\) 表示。
\end{enumerate}

\subsection{八、风险与提醒}\label{ux516bux98ceux9669ux4e0eux63d0ux9192}

\begin{itemize}
\tightlist
\item
  ⚠️ Kaggle 双 T4 会话\textbf{不稳定}（可能只给一张）：预检时确认 \texttt{torch.cuda.device\_count()==2}。
\item
  ⚠️ 不要在本课题纠缠''调到最快''------\textbf{目标是理解扩展效率的成因}，配额要留给课题08。
\end{itemize}

\end{document}
